\documentclass[aspectratio=169]{beamer}

\usetheme{Madrid}
\usecolortheme{default}

\usepackage[utf8]{inputenc}
\usepackage[T1]{fontenc}
\usepackage[brazil]{babel}
\usepackage{amsmath}
\usepackage{amsfonts}
\usepackage{graphicx}
\usepackage{booktabs}
\usepackage{xcolor}

\usepackage{algorithm}
\usepackage{algpseudocode}

\title[DDPM para Geração de Faces]{Modelo de Difusão DDPM para Geração de Faces}
\subtitle{Arquitetura, Treinamento e Amostragem}
\author{Lorran de Araújo Durães Soares}
\institute{Tópicos Avançados em Redes Neurais Profundas}
\date{\today}

\begin{document}

\begin{frame}
    \titlepage
\end{frame}

\begin{frame}{Objetivo}
    \begin{itemize}
        \item Treinar um modelo generativo capaz de produzir imagens de faces a partir de ruído aleatório.
        \item O modelo usado é um DDPM: \textit{Denoising Diffusion Probabilistic Model} \cite{ho2020ddpm}.
        \item A rede neural principal é uma UNet convolucional condicionada pelo passo temporal de difusão \cite{ronneberger2015unet}.
        \item A geração é feita por remoção progressiva de ruído, do passo $T$ até o passo $0$.
    \end{itemize}
\end{frame}

\begin{frame}{Fluxo Geral do Projeto}
    \begin{enumerate}
        \item Normalizar as imagens para o intervalo $[-1, 1]$.
        \item Durante o treinamento, adicionar ruído gaussiano às imagens.
        \item Treinar uma UNet para prever o ruído adicionado.
        \item Durante a geração, começar de ruído puro.
        \item Aplicar o processo reverso até obter uma imagem sintética.
    \end{enumerate}
\end{frame}

\begin{frame}{Dataset e Pré-processamento}
    \begin{itemize}
        \item Dataset usado: subconjunto do UTKFace \cite{zhang2017utkface}.
        \item Filtro aplicado: faces femininas entre 20 e 40 anos.
        \item Imagens convertidas para RGB e redimensionadas para $128 \times 128$.
        \item Normalização:
        \[
            x \in [0, 1] \rightarrow x \in [-1, 1]
        \]
        \item Aumento de dados: \texttt{RandomHorizontalFlip} com probabilidade 0.5.
    \end{itemize}
\end{frame}

\begin{frame}{Ideia Central do DDPM}
    \begin{block}{Processo direto}
        Uma imagem real é corrompida progressivamente com ruído gaussiano.
    \end{block}

    \begin{block}{Processo reverso}
        O modelo aprende a remover esse ruído passo a passo.
    \end{block}

    \begin{itemize}
        \item No treinamento, conhecemos a imagem limpa e o ruído aplicado.
        \item A rede aprende a prever o ruído presente em uma imagem ruidosa.
        \item Na geração, começamos de ruído puro e usamos a rede para reverter o processo.
    \end{itemize}
\end{frame}

\begin{frame}{Processo Direto: Adição de Ruído}
    O scheduler define uma sequência linear de variâncias $\beta_t$ \cite{ho2020ddpm}.

    \[
        \alpha_t = 1 - \beta_t
    \]

    \[
        \bar{\alpha}_t = \prod_{s=1}^{t} \alpha_s
    \]

    A imagem ruidosa em um passo $t$ é obtida por:

    \[
        x_t = \sqrt{\bar{\alpha}_t} x_0 + \sqrt{1 - \bar{\alpha}_t} \epsilon
    \]

    onde $\epsilon \sim \mathcal{N}(0, I)$.
\end{frame}

\begin{frame}{Objetivo de Treinamento}
    A rede neural recebe:

    \[
        (x_t, t)
    \]

    e prediz:

    \[
        \epsilon_\theta(x_t, t)
    \]

    A função de perda é \cite{ho2020ddpm}:

    \[
        \mathcal{L} = \mathbb{E}_{x_0, t, \epsilon}\left[\|\epsilon - \epsilon_\theta(x_t, t)\|^2\right]
    \]

    Ou seja, o modelo não prediz diretamente a imagem final; ele prediz o ruído.
\end{frame}

\begin{frame}{Arquitetura Geral da UNet I}

    \begin{figure}
        \centering
        \includegraphics[width=0.85\textwidth]{imgs/arquitetura.png}
    \end{figure}

\end{frame}

\begin{frame}{Arquitetura Geral da UNet II}
    \begin{itemize}
        \item Entrada: imagem ruidosa $x_t$ com forma $[B, 3, 128, 128]$.
        \item Condicionamento: timestep $t$ via embedding temporal.
        \item Saída: ruído previsto com a mesma forma da imagem.
        \item Estrutura encoder-decoder com conexões de atalho \cite{ronneberger2015unet}.
        \item Encoder captura contexto global.
        \item Decoder reconstrói a resolução espacial.
        \item Skip connections preservam detalhes locais.
    \end{itemize}
\end{frame}

\begin{frame}{Escalas da UNet}
    \begin{table}
        \centering
        \begin{tabular}{ccc}
            \toprule
            Etapa & Resolução & Canais \\
            \midrule
            Entrada & $128 \times 128$ & 3 \\
            Nível 1 & $128 \times 128$ & 64 \\
            Nível 2 & $64 \times 64$ & 128 \\
            Nível 3 & $32 \times 32$ & 256 \\
            Nível 4 & $16 \times 16$ & 512 \\
            Bottleneck & $8 \times 8$ & 512 \\
            Saída & $128 \times 128$ & 3 \\
            \bottomrule
        \end{tabular}
    \end{table}

    Com \texttt{base\_channels = 64}, o modelo possui aproximadamente 38.49 milhões de parâmetros.
\end{frame}

\begin{frame}{Bloco Residual}
    Cada \texttt{ResBlock} contém:

    \begin{itemize}
        \item GroupNorm \cite{wu2018groupnorm}.
        \item Ativação SiLU.
        \item Convolução $3 \times 3$.
        \item Injeção do embedding temporal.
        \item Nova sequência GroupNorm, SiLU e convolução $3 \times 3$.
        \item Conexão residual entre entrada e saída \cite{he2016resnet}.
    \end{itemize}

    \[
        y = F(x, t) + x
    \]

    Se o número de canais muda, usa-se uma convolução $1 \times 1$ no caminho residual.
\end{frame}

\begin{frame}{Embedding Temporal}
    O modelo precisa saber o nível de ruído de cada amostra.

    \begin{itemize}
        \item O timestep $t$ é codificado com senos e cossenos em múltiplas frequências.
        \item Essa estratégia é inspirada nos embeddings posicionais dos Transformers \cite{vaswani2017attention}.
        \item Depois, um MLP projeta o embedding para dimensão 256.
        \item O embedding é somado aos mapas de características dentro dos ResBlocks.
    \end{itemize}

    Injeção temporal no bloco:

    \[
        h = h + \text{MLP}(t)[:, :, None, None]
    \]
\end{frame}

\begin{frame}{Self-Attention}
    \begin{itemize}
        \item A atenção permite que regiões distantes da imagem interajam.
        \item Isso é útil para faces, pois olhos, boca, cabelo e contorno precisam ser globalmente coerentes.
        \item O modelo usa atenção nas resoluções menores: $16 \times 16$ e $8 \times 8$.
        \item Aplicar atenção em $128 \times 128$ seria muito custoso em memória.
    \end{itemize}

    Atenção por produto escalar \cite{vaswani2017attention}:

    \[
        \text{Attention}(Q, K, V) = \text{softmax}\left(\frac{QK^T}{\sqrt{C}}\right)V
    \]
\end{frame}

\begin{frame}{Treinamento do DDPM}
    \begin{figure}
        \centering
        \includegraphics[width=0.55\textwidth]{imgs/algorithm1.png}
    \end{figure}
\end{frame}

\begin{frame}{Loop de Treinamento}
    Para cada batch de imagens reais $x_0$:

    \begin{enumerate}
        \item Sorteia-se um timestep $t$ para cada imagem.
        \item Sorteia-se ruído gaussiano $\epsilon$.
        \item Calcula-se $x_t$ usando o scheduler.
        \item A UNet prediz $\epsilon_\theta(x_t, t)$.
        \item Calcula-se a loss MSE.
        \item Executa-se backpropagation.
        \item Aplica-se gradient clipping.
        \item Atualiza-se o modelo com AdamW \cite{loshchilov2019adamw}.
        \item Atualiza-se a EMA dos pesos.
    \end{enumerate}
\end{frame}

\begin{frame}{Hiperparâmetros Principais}
    \begin{table}
        \centering
        \begin{tabular}{ll}
            \toprule
            Hiperparâmetro & Valor atual \\
            \midrule
            Tamanho da imagem & 128 \\
            Batch size & 16 \\
            Épocas & 200 \\
            Learning rate & 0.0002 \\
            Timesteps & 1000 \\
            Base channels & 64 \\
            EMA decay & 0.999 \\
            Gradient clipping & 1.0 \\
            Augmentation & Horizontal flip \\
            \bottomrule
        \end{tabular}
    \end{table}
\end{frame}

\begin{frame}{EMA: Média Móvel Exponencial}
    A EMA mantém uma segunda cópia do modelo, suavizando os pesos ao longo do treinamento.

    \[
        \theta_{EMA} \leftarrow \lambda \theta_{EMA} + (1 - \lambda)\theta
    \]

    No projeto:

    \[
        \lambda = 0.999
    \]

    Logo:

    \[
        \theta_{EMA} \leftarrow 0.999\theta_{EMA} + 0.001\theta
    \]
\end{frame}

\begin{frame}{Por que Usar EMA?}
    \begin{itemize}
        \item Os pesos do modelo principal oscilam durante o treinamento.
        \item A oscilação vem dos batches aleatórios, timesteps aleatórios e ruído gaussiano.
        \item A EMA cria uma versão temporalmente suavizada dos pesos \cite{nichol2021improved,dhariwal2021diffusion}.
        \item Em modelos generativos, a EMA costuma gerar imagens mais estáveis e coerentes.
        \item A EMA não altera a loss: o modelo principal continua treinando com MSE.
    \end{itemize}
\end{frame}

\begin{frame}{Processo de Geração I}
    \begin{figure}
        \centering
        \includegraphics[width=0.5\textwidth]{imgs/algorithm2.png}
    \end{figure}
\end{frame}

\begin{frame}{Processo de Geração II}
    A geração começa com ruído puro:

    \[
        x_T \sim \mathcal{N}(0, I)
    \]

    Depois, para cada timestep $t$ de $T-1$ até $0$:

    \begin{enumerate}
        \item A UNet prediz o ruído presente em $x_t$.
        \item Calcula-se a média da distribuição reversa.
        \item Se $t > 0$, adiciona-se ruído com variância posterior.
        \item Se $t = 0$, a média é usada como imagem final.
    \end{enumerate}
\end{frame}

\begin{frame}{Exemplos de imagens geradas}
    \begin{figure}
        \centering
        \includegraphics[width=0.85\textwidth]{imgs/generateds.png}
    \end{figure}
\end{frame}

\begin{frame}{Visualização do processo de Denoising}
    \begin{figure}
        \centering
        \includegraphics[width=0.85\textwidth]{imgs/denoising.png}
    \end{figure}
\end{frame}

\begin{frame}{Limitações Atuais}
    \begin{itemize}
        \item A geração DDPM com 1000 passos é lenta.
        \item O modelo é incondicional: não controla idade, gênero ou etnia explicitamente.
        \item Não há métrica quantitativa como FID implementada.
        \item O scheduler de ruído é linear.
        \item O conjunto filtrado tem tamanho limitado para geração de faces.
    \end{itemize}
\end{frame}

\begin{frame}{Possíveis Melhorias}
    \begin{itemize}
        \item Implementar DDIM para amostragem mais rápida \cite{song2021ddim}.
        \item Testar cosine noise schedule.
        \item Usar mixed precision para acelerar o treinamento.
        \item Adicionar condicionamento por atributos do UTKFace.
        \item Medir qualidade com FID.
        \item Treinar por mais tempo ou usar dataset maior.
    \end{itemize}
\end{frame}

\begin{frame}{Resumo}
    \begin{itemize}
        \item O DDPM aprende a inverter um processo de adição gradual de ruído.
        \item A UNet prediz o ruído em cada timestep.
        \item A arquitetura usa múltiplas escalas, ResBlocks, skip connections e atenção \cite{ronneberger2015unet,he2016resnet,vaswani2017attention}.
        \item O treinamento usa MSE, AdamW, gradient clipping e EMA.
        \item Na geração, a versão EMA do modelo é preferida por produzir amostras mais estáveis.
    \end{itemize}
\end{frame}

\begin{frame}{Referências}
    \footnotesize
    \begin{thebibliography}{99}
        \bibitem{ho2020ddpm} Ho, J.; Jain, A.; Abbeel, P. \textit{Denoising Diffusion Probabilistic Models}. NeurIPS, 2020.
        \bibitem{ronneberger2015unet} Ronneberger, O.; Fischer, P.; Brox, T. \textit{U-Net: Convolutional Networks for Biomedical Image Segmentation}. MICCAI, 2015.
        \bibitem{he2016resnet} He, K.; Zhang, X.; Ren, S.; Sun, J. \textit{Deep Residual Learning for Image Recognition}. CVPR, 2016.
        \bibitem{vaswani2017attention} Vaswani, A. et al. \textit{Attention Is All You Need}. NeurIPS, 2017.
        \bibitem{wu2018groupnorm} Wu, Y.; He, K. \textit{Group Normalization}. ECCV, 2018.
        \bibitem{loshchilov2019adamw} Loshchilov, I.; Hutter, F. \textit{Decoupled Weight Decay Regularization}. ICLR, 2019.
        \bibitem{nichol2021improved} Nichol, A. Q.; Dhariwal, P. \textit{Improved Denoising Diffusion Probabilistic Models}. ICML, 2021.
        \bibitem{dhariwal2021diffusion} Dhariwal, P.; Nichol, A. Q. \textit{Diffusion Models Beat GANs on Image Synthesis}. NeurIPS, 2021.
        \bibitem{song2021ddim} Song, J.; Meng, C.; Ermon, S. \textit{Denoising Diffusion Implicit Models}. ICLR, 2021.
        \bibitem{zhang2017utkface} Zhang, Z.; Song, Y.; Qi, H. \textit{Age Progression/Regression by Conditional Adversarial Autoencoder}. CVPR, 2017.
    \end{thebibliography}
\end{frame}

\end{document}
